% GENERATED FILE. Edit canonical lessons, then run npm run generate:books.
% canonical-source-hash: fnv1a64:d136eec08740c9d4
% canonical-lessons: KA-C342-muttige-haku, KA-C342-kandaka, KA-C342-suttuvari, KA-C342-saranagu, KA-C342-phirangi

\chapter{The Siege of the Fort}
\label{ch:ka-the-siege-of-the-fort}
\hlchaptermodality{342}

\begin{chapteropening}
\textbf{By the end of this chapter:} \emph{I can say to besiege, a moat, to surround, to surrender and a cannon.}

Complete the last lesson of chapter 342: I can say to besiege, a moat, to surround, to surrender and a cannon.
\end{chapteropening}

\section[\texorpdfstring{\kn{ಮುತ್ತಿಗೆ} \kn{ಹಾಕು} (muttige hāku)}{ಮುತ್ತಿಗೆ ಹಾಕು (muttige hāku)}]{\kn{ಮುತ್ತಿಗೆ} \kn{ಹಾಕು} (muttige hāku) — to besiege}

\label{lesson:KA-C342-muttige-haku}



\begin{quote}
Before the new one: say the Kannada for a family line, then the Kannada for a monarch.
\end{quote}

\subsection*{You'll want to know: \kn{ಮುತ್ತಿಗೆ} \kn{ಹಾಕು}}
\textbf{\kn{ಮುತ್ತಿಗೆ} \kn{ಹಾಕು}} — \emph{muttige hāku} — \textquotedblleft{}to besiege\textquotedblright{}.

\begin{quote}
\kn{ಆ} \kn{ದೊಡ್ಡ} \kn{ಸೈನ್ಯ} \kn{ಕೋಟೆಗೆ} \kn{ಮುತ್ತಿಗೆ} \kn{ಹಾಕಿತು}.
\end{quote}

\begin{quote}
\emph{ā doḍḍa sainya kōṭege muttige hākitu.}
\end{quote}

\begin{quote}
That great army besieged the fort.
\end{quote}

Past (I …): \kn{ಮುತ್ತಿಗೆ} \kn{ಹಾಕಿದೆ}, \emph{muttige hākide}.

What is besieged takes -\kn{ಗೆ}. Today people also \kn{ಮುತ್ತಿಗೆ} \kn{ಹಾಕು} an office in protest.

\begin{cousinweb}
\textbf{\kn{ಮುತ್ತಿಗೆ} \kn{ಹಾಕು}} (\emph{muttige hāku}) is native Dravidian: \kn{ಮುತ್ತಿಗೆ}, \textquotedblleft{}a siege\textquotedblright{}, belongs with Kannada \emph{muttu}, \textquotedblleft{}to enclose\textquotedblright{}, and Tamil \emph{muṟṟu}, \textquotedblleft{}to surround, besiege\textquotedblright{} (DEDR 5018), with \kn{ಹಾಕು}, \textquotedblleft{}to put\textquotedblright{}.
\end{cousinweb}

\begin{grammarlens}[title={What you've built}]
One more word for telling what happened.
\end{grammarlens}

\subsection*{Your turn}
\begin{itemize}
\raggedright
  \item \emph{Say it:} \emph{muttige hāku}
  \item \emph{Say it:} \emph{muttige hāku}, once more
  \item \emph{Say it:} \emph{ā doḍḍa sainya kōṭege muttige hākitu.} — That great army besieged the fort.
  \item \emph{Recall:} say the Kannada for a family line, then the Kannada for a monarch, then say \emph{muttige hāku} again
\end{itemize}

\begin{tcolorbox}[breakable,colback=teal!4,colframe=teal!35!black,title={Before you move on}]
What does \kn{ಮುತ್ತಿಗೆ} \kn{ಹಾಕು} mean? (\textquotedblleft{}To besiege\textquotedblright{}.) Say it once more.
\end{tcolorbox}
\section[\texorpdfstring{\kn{ಕಂದಕ} (kandaka)}{ಕಂದಕ (kandaka)}]{\kn{ಕಂದಕ} (kandaka) — a moat; a trench}

\label{lesson:KA-C342-kandaka}



\begin{quote}
Before the new one: say the Kannada for a monarch, then the Kannada for to besiege.
\end{quote}

\subsection*{You'll want to know: \kn{ಕಂದಕ}}
\textbf{\kn{ಕಂದಕ}} — \emph{kandaka} — \textquotedblleft{}a moat; a trench\textquotedblright{}.

\begin{quote}
\kn{ಕೋಟೆಯ} \kn{ಸುತ್ತ} \kn{ಆಳವಾದ} \kn{ಕಂದಕ} \kn{ಇತ್ತು}.
\end{quote}

\begin{quote}
\emph{kōṭeya sutta āḷavāda kandaka ittu.}
\end{quote}

\begin{quote}
There was a deep moat around the fort.
\end{quote}

A ditch dug for defence, often full of water.

\begin{cousinweb}
\textbf{\kn{ಕಂದಕ}} (\emph{kandaka}) is usually traced to Arabic \emph{khandaq}, \textquotedblleft{}a ditch, a trench\textquotedblright{}, which came into the Indian languages through Persian.
\end{cousinweb}

\begin{grammarlens}[title={What you've built}]
One more word for telling what happened.
\end{grammarlens}

\subsection*{Your turn}
\begin{itemize}
\raggedright
  \item \emph{Say it:} \emph{kandaka}
  \item \emph{Say it:} \emph{kandaka}, once more
  \item \emph{Say it:} \emph{kōṭeya sutta āḷavāda kandaka ittu.} — There was a deep moat around the fort.
  \item \emph{Recall:} say the Kannada for a monarch, then the Kannada for to besiege, then say \emph{kandaka} again
\end{itemize}

\begin{tcolorbox}[breakable,colback=teal!4,colframe=teal!35!black,title={Before you move on}]
What does \kn{ಕಂದಕ} mean? (\textquotedblleft{}A moat; a trench\textquotedblright{}.) Say it once more.
\end{tcolorbox}
\section[\texorpdfstring{\kn{ಸುತ್ತುವರಿ} (suttuvari)}{ಸುತ್ತುವರಿ (suttuvari)}]{\kn{ಸುತ್ತುವರಿ} (suttuvari) — to surround}

\label{lesson:KA-C342-suttuvari}



\begin{quote}
Before the new one: say the Kannada for to besiege, then the Kannada for a moat.
\end{quote}

\subsection*{You'll want to know: \kn{ಸುತ್ತುವರಿ}}
\textbf{\kn{ಸುತ್ತುವರಿ}} — \emph{suttuvari} — \textquotedblleft{}to surround\textquotedblright{}.

\begin{quote}
\kn{ಪೊಲೀಸರು} \kn{ಆ} \kn{ಮನೆಯನ್ನು} \kn{ಸುತ್ತುವರಿದರು}.
\end{quote}

\begin{quote}
\emph{polīsaru ā maneyannu suttuvaridaru.}
\end{quote}

\begin{quote}
The police surrounded that house.
\end{quote}

Past (I …): \kn{ಸುತ್ತುವರಿದೆ}, \emph{suttuvaride}.

What is surrounded takes -\kn{ಅನ್ನು}.

\begin{cousinweb}
\textbf{\kn{ಸುತ್ತುವರಿ}} (\emph{suttuvari}) is built on native \emph{suttu}, \textquotedblleft{}to surround, go round\textquotedblright{} (DEDR 2715), as in \kn{ಸುತ್ತ}, \textquotedblleft{}around\textquotedblright{}; the origin of the second part is not settled.
\end{cousinweb}

\begin{grammarlens}[title={What you've built}]
One more word for telling what happened.
\end{grammarlens}

\subsection*{Your turn}
\begin{itemize}
\raggedright
  \item \emph{Say it:} \emph{suttuvari}
  \item \emph{Say it:} \emph{suttuvari}, once more
  \item \emph{Say it:} \emph{polīsaru ā maneyannu suttuvaridaru.} — The police surrounded that house.
  \item \emph{Recall:} say the Kannada for to besiege, then the Kannada for a moat, then say \emph{suttuvari} again
\end{itemize}

\begin{tcolorbox}[breakable,colback=teal!4,colframe=teal!35!black,title={Before you move on}]
What does \kn{ಸುತ್ತುವರಿ} mean? (\textquotedblleft{}To surround\textquotedblright{}.) Say it once more.
\end{tcolorbox}
\section[\texorpdfstring{\kn{ಶರಣಾಗು} (śaraṇāgu)}{ಶರಣಾಗು (śaraṇāgu)}]{\kn{ಶರಣಾಗು} (śaraṇāgu) — to surrender}

\label{lesson:KA-C342-saranagu}



\begin{quote}
Before the new one: say the Kannada for a moat, then the Kannada for to surround.
\end{quote}

\subsection*{You'll want to know: \kn{ಶರಣಾಗು}}
\textbf{\kn{ಶರಣಾಗು}} — \emph{śaraṇāgu} — \textquotedblleft{}to surrender\textquotedblright{}.

\begin{quote}
\kn{ಊಟ} \kn{ಇಲ್ಲದೆ} \kn{ಕೋಟೆಯ} \kn{ಸೈನಿಕರು} \kn{ಶರಣಾದರು}.
\end{quote}

\begin{quote}
\emph{ūṭa illade kōṭeya sainikaru śaraṇādaru.}
\end{quote}

\begin{quote}
With no food, the soldiers in the fort surrendered.
\end{quote}

Past (I …): \kn{ಶರಣಾದೆ}, \emph{śaraṇāde}.

The one surrendered to takes -\kn{ಗೆ}: \kn{ಪೊಲೀಸರಿಗೆ} \kn{ಶರಣಾಗು}.

\begin{cousinweb}
\textbf{\kn{ಶರಣಾಗು}} (\emph{śaraṇāgu}) is built from Sanskrit \emph{śaraṇa}, \textquotedblleft{}protecting, shelter\textquotedblright{} (CDIAL 12326), and native \emph{āgu}, \textquotedblleft{}to become\textquotedblright{} (DEDR 333): to seek shelter with the stronger side.
\end{cousinweb}

\begin{grammarlens}[title={What you've built}]
One more word for telling what happened.
\end{grammarlens}

\subsection*{Your turn}
\begin{itemize}
\raggedright
  \item \emph{Say it:} \emph{śaraṇāgu}
  \item \emph{Say it:} \emph{śaraṇāgu}, once more
  \item \emph{Say it:} \emph{ūṭa illade kōṭeya sainikaru śaraṇādaru.} — With no food, the soldiers in the fort surrendered.
  \item \emph{Recall:} say the Kannada for a moat, then the Kannada for to surround, then say \emph{śaraṇāgu} again
\end{itemize}

\begin{tcolorbox}[breakable,colback=teal!4,colframe=teal!35!black,title={Before you move on}]
What does \kn{ಶರಣಾಗು} mean? (\textquotedblleft{}To surrender\textquotedblright{}.) Say it once more.
\end{tcolorbox}
\section[\texorpdfstring{\kn{ಫಿರಂಗಿ} (phiraṅgi)}{ಫಿರಂಗಿ (phiraṅgi)}]{\kn{ಫಿರಂಗಿ} (phiraṅgi) — a cannon}

\label{lesson:KA-C342-phirangi}



\begin{quote}
Before the new one: say the Kannada for to surround, then the Kannada for to surrender.
\end{quote}

\subsection*{You'll want to know: \kn{ಫಿರಂಗಿ}}
\textbf{\kn{ಫಿರಂಗಿ}} — \emph{phiraṅgi} — \textquotedblleft{}a cannon\textquotedblright{}.

\begin{quote}
\kn{ಕೋಟೆಯ} \kn{ಗೋಡೆಯ} \kn{ಮೇಲೆ} \kn{ಹಳೆಯ} \kn{ಫಿರಂಗಿ} \kn{ಇದೆ}.
\end{quote}

\begin{quote}
\emph{kōṭeya gōḍeya mēle haḷeya phiraṅgi ide.}
\end{quote}

\begin{quote}
There is an old cannon on the fort wall.
\end{quote}

Plural: \kn{ಫಿರಂಗಿಗಳು}, \emph{phiraṅgigaḷu}.

Many old forts still show their \kn{ಫಿರಂಗಿಗಳು}.

\begin{cousinweb}
\textbf{\kn{ಫಿರಂಗಿ}} (\emph{phiraṅgi}) is a loan from Persian \emph{firangī}, \textquotedblleft{}a Frank, a European\textquotedblright{}: the guns came with European traders and soldiers.
\end{cousinweb}

\begin{grammarlens}[title={What you've built}]
One more word for telling what happened.
\end{grammarlens}

\subsection*{Your turn}
\begin{itemize}
\raggedright
  \item \emph{Say it:} \emph{phiraṅgi}
  \item \emph{Say it:} \emph{phiraṅgi}, once more
  \item \emph{Say it:} \emph{kōṭeya gōḍeya mēle haḷeya phiraṅgi ide.} — There is an old cannon on the fort wall.
  \item \emph{Recall:} say the Kannada for to surround, then the Kannada for to surrender, then say \emph{phiraṅgi} again
\end{itemize}

\begin{tcolorbox}[breakable,colback=teal!4,colframe=teal!35!black,title={Before you move on}]
What does \kn{ಫಿರಂಗಿ} mean? (\textquotedblleft{}A cannon\textquotedblright{}.) Say it once more.
\end{tcolorbox}
